\chapter{環論}
この章では環論の問題を扱う。環論は代数幾何学, 整数論を展開するために基礎となる分野であり, なかなか手強いものである。参考書としては\cite{atiyah}が有名である。なお, スキーム論の対策はこのpdfの専門のほうに述べる. \par
唐突だが, 次の問題に挑戦してみよう。次の問題は, 設問が5つあるが, だんだんレベルが上がっていくような構成が見事な問題である。
\begin{prac} (神戸大R1-A) \label{prac:kobeR1_A}
$p$を素数とする。
\[R = \set{\dfrac{m}{n}\in \mb{Q}\mid m,n\in \mb{Z}, p\not | n}\]
とおく。次の問いに答えよ。
\ben
\item $R$が$\mb{Q}$の部分環であることを示せ。
\item $p$が生成する$R$のイデアルが極大であることを示せ。
\item (2)の極大イデアル以外に$R$が極大イデアルを持たないことを示せ。
\item $R$係数$1$変数多項式環$R[X]$の$X^2+1$で生成されるイデアルを$I$とおく。$I$が素イデアルであることを示せ。
\item (4)のイデアル$I$を含む$R[X]$の極大イデアルを$J$とおく。$J$が$p$を含むことを示せ。
\een
\end{prac}
(5)は単純に見えるが, 可換環論の知識を用いるとテクニカルに導ける。(4)までなら速く解けて, かつこの(5)ができる人は可換環論に一定の親しみがあると思われる。(4)までできれば基礎力はある。ひとまず(4)までを自力でできるようになるとよい。(5)が分からなかった人は, \cite{atiyah}の5章(整従属)のあたりを復習してみよう。
(5)の答えのみ脚注\footnote{
$R$代数の同型$R[X]/(X^2+1)\cong R[\sqrt{-1}]$がある。包含写像$i:R\to R[i]$は単射整準同型であるから, $i$が誘導する$\spec{i}: \spec{R[i]}\to \spec{R}$は極大イデアルを極大イデアルに移す。なぜなら, $R/i^{-1}(J)\to R[i]/J$は整域の間の単射整準同型で,$R[i]/J$が体であることと$R/i^{-1}(J)$が体であることが同値であるためである。よって$i^{-1}(J) = J\cap R$は$R$の極大イデアルだから$J\cap R = (p)$である。ゆえに$p\in J$を得る。\qed

}に与えておく。\par





重要な環と簡単な性質について復習する。\tf{以降ことわりの無い限り環は単位的可換環を指す。}
